\section{农业机械中的长幅摆线}
\markboth{20. 农业机械中的长幅摆线}{20. 农业机械中的长幅摆线}

很多农业机械在工作时，都有一些部件在绕着自己的轴旋转，这些部件又随着整个机具一起向前作匀速直线运动。根据第 19 节的讨论，在这些一边旋转一边前进的部件上，每个点都在空中画出一条摆线或变幅摆线。

农业机械里经常用到长幅摆线，而且主要应用它的绕扣部分。例如，图 20.1 是卧式旋耕机的工作原理示意图\footnote{重排校订：原书图 20.1 中刀片编号“1”位置不清，且四条长幅摆线轨迹起点截断参差不齐；本版重绘时将四条轨迹统一自各刀尖（$t=0$）同步出发，并清晰标出刀片 1～4。}。从图中可以看出，卧式旋耕机的每把刀片画出一条长幅摆线，它的绕扣部分很大。调整旋转轴的高度，可以使刀片在绕扣最宽的地方入土，翻松绕扣下半截的泥土后，再露出地面。四把刀片画出的四条长幅摆线顺次错开，绕扣部分互相重迭，实现了连续翻耕土地的动作。在这种机具上只有利用长幅摆线才行，如果刀尖的运动轨迹不是长幅摆线，而是普通摆线或短幅摆线，那么就不能进行翻土的动作。

\begin{center}
  \begin{tikzpicture}[scale=1.02, >=Stealth, font=\footnotesize]
    \def\r{0.26}
    \def\a{0.64}

    % 水平轴线与中心十字线
    \draw[thin] ({-\a - 0.14}, 0) -- ({3*pi*\r + \a + 0.20}, 0);
    \draw[thin] (0, {-\a - 0.14}) -- (0, {\a + 0.14});

    % 旋耕机圆盘虚线外圈
    \draw[dashed, mutedgray] (0, 0) circle (\a);

    % 四把弯刀片（1 上，2 左，3 下，4 右）
    \foreach \ang in {90, 180, 270, 0} {
      \begin{scope}[rotate=\ang]
        \draw[thick] (0, 0) .. controls ({0.35*\a}, {0.15*\a}) and ({0.65*\a}, {0.35*\a}) .. (\a, 0);
        \draw[thick] (0, 0) .. controls ({0.35*\a}, {-0.05*\a}) and ({0.72*\a}, {0.15*\a}) .. (\a, 0);
      \end{scope}
    }

    % 刀盘中心顺时针旋转弧形箭头（匹配向右行进速度 u）
    \draw[-{Stealth[length=3.6pt, width=2.6pt]}, line width=0.65pt]
      ({0.26*cos(76)}, {0.26*sin(76)}) arc[start angle=76, end angle=14, radius=0.26];

    % 前进方向 u 与刀片标注
    \draw[->] (-0.88, {\a + 0.28}) -- (-0.44, {\a + 0.28});
    \node[above, inner sep=1pt] at (-0.66, {\a + 0.29}) {$u$};
    \node[above] at (0, {\a + 0.16}) {刀片};

    % 四把刀片的长幅摆线轨迹（均从 t=0 各刀尖出发，右端统一整齐收束于最后绕扣 1 对称轴 x=3*pi*r 附近）
    % 刀片 1（90°上）：形成绕扣 1(pi), 1(3pi)，于自交点闭合
    \draw[thick, curveblue, smooth, domain=0:{3.672*pi}, samples=130]
      plot ({\r*\x + \a*sin(\x r)}, {\a*cos(\x r)});
    % 刀片 2（180°左）：形成绕扣 2(1.5pi)，于 x=3*pi*r 处与刀片 4 对称交汇
    \draw[thick, curveblue, smooth, domain=0:{2.648*pi}, samples=100]
      plot ({\r*\x + \a*sin((\x - 0.5*pi) r)}, {\a*cos((\x - 0.5*pi) r)});
    % 刀片 3（270°下）：形成绕扣 3(2pi)，于 x=3*pi*r 处到达波峰收束
    \draw[thick, curveblue, smooth, domain=0:{3.00*pi}, samples=110]
      plot ({\r*\x + \a*sin((\x - pi) r)}, {\a*cos((\x - pi) r)});
    % 刀片 4（0°右）：形成绕扣 4(0.5pi), 4(2.5pi)，于 x=3*pi*r 处与刀片 2 对称交汇
    \draw[thick, curveblue, smooth, domain=0:{3.352*pi}, samples=120]
      plot ({\r*\x + \a*sin((\x - 1.5*pi) r)}, {\a*cos((\x - 1.5*pi) r)});

    % 刀片编号 1~4 置于所有蓝色轨迹之后，加白底遮罩放于刀盘内侧干净空白处
    \node[fill=white, inner sep=0.8pt] at (0.16, 0.48) {1};
    \node[fill=white, inner sep=0.8pt] at (-0.46, 0.16) {2};
    \node[fill=white, inner sep=0.8pt] at (-0.16, -0.46) {3};
    \node[fill=white, inner sep=0.8pt] at (0.45, -0.18) {4};

    % 底部绕扣编号 4 1 2 3 4 1（与全部 6 个绕扣一一严格对应）
    \node[below] at ({\r*0.5*pi}, {-\a - 0.03}) {4};
    \node[below] at ({\r*1.0*pi}, {-\a - 0.03}) {1};
    \node[below] at ({\r*1.5*pi}, {-\a - 0.03}) {2};
    \node[below] at ({\r*2.0*pi}, {-\a - 0.03}) {3};
    \node[below] at ({\r*2.5*pi}, {-\a - 0.03}) {4};
    \node[below] at ({\r*3.0*pi}, {-\a - 0.03}) {1};
  \end{tikzpicture}
  \captionof*{figure}{图 20.1}
\end{center}

又如联合收割机的拨禾板，一面跟随联合收割机前进，一面绕自己的轴旋转。拨禾板要能沿竖直方向插入禾丛，把禾杆引向切割器，并把割断的禾杆向后推送，所以拨禾板的外端也应该描绘长幅摆线，并且在绕扣部分的下半截曲线弧上进行工作。
